% NLP MAI SS2026 – Project Proposal
% Compile (recommended):
%   pdflatex proposal.tex
% or
%   latexmk -pdf proposal.tex

\documentclass[11pt,a4paper]{article}

\usepackage[a4paper,margin=2.2cm]{geometry}
\usepackage[T1]{fontenc}
\usepackage[utf8]{inputenc}
\usepackage{lmodern}
\usepackage{microtype}
\usepackage{enumitem}
\usepackage{hyperref}
\hypersetup{colorlinks=true,linkcolor=blue,urlcolor=blue,citecolor=blue}

\title{\vspace{-1.0cm}
\textbf{Project Proposal (NLP MAI SS2026)}\\
\vspace{0.2cm}
\textbf{``Who Hunts Who?'' Comparing Three Wikipedia-based Knowledge Graph Construction Methods for Predator--Prey Relations}
\vspace{-0.2cm}}
\author{Group Members: \; Zoltan Bek,  Armin Lohse, Nino Reil, Andreas Wallnöfer}

\begin{document}
\maketitle
\vspace{-0.4cm}

\section*{1. Motivation \& Problem Definition}
Knowledge graphs (KGs) represent real-world \textit{entities} and their \textit{relationships} as a graph. In practice, building a KG from text is not unique: different construction methods can yield different edges and graph structures, even on the same corpus. This project uses a small and intuitive domain---\textbf{predator--prey relations in animals}---to compare three KG construction approaches on Wikipedia data and assess which approach best captures the directed relation \textbf{predator $\rightarrow$ prey}.

\section*{2. Task}
We build and evaluate a small knowledge graph where:
\begin{itemize}[leftmargin=*,nosep]
  \item Nodes represent animal entities.
  \item Directed edges represent candidate predator--prey relations (\textit{``A eats/hunts B''}).
\end{itemize}
All work is delivered as a \textbf{single Jupyter notebook} (no additional project structure required).

\section*{3. Research Question}
\textbf{How do three different knowledge graph construction methods differ in accuracy (Precision/Recall/F1) when extracting predator$\rightarrow$prey relations from Wikipedia, and what are their runtime/graph-structure trade-offs?}

\section*{4. Dataset (small, $\sim$20 pages)}
We use a controlled Wikipedia corpus of approximately \textbf{20 general animal pages} chosen to form a terrestrial food web where multi-step chains are possible (e.g., \textit{A eats B, B eats C}).

\textbf{Predators / higher trophic:} Hawk, Eagle, Owl, Fox, Wolf, Snake \\
\textbf{Mid-level (predator + prey):} Frog, Lizard, Spider, Bird, Shrew, Bat \\
\textbf{Lower trophic / common prey:} Mouse, Rabbit, Squirrel, Grasshopper, Cricket, Caterpillar, Fly, Beetle

This selection supports chain examples such as \textit{Hawk $\rightarrow$ Snake $\rightarrow$ Frog $\rightarrow$ Grasshopper}. We start with general pages for simplicity and breadth; during implementation we may refine selected nodes to \textbf{specific species pages} to obtain clearer diet statements for gold-standard labeling.

\section*{5. Methods (3 KG Construction Approaches)}
We construct three graphs from the same dataset; each produces a set of directed candidate edges predator$\rightarrow$prey.

\begin{itemize}[leftmargin=*,itemsep=0.25em]
  \item \textbf{Method A: Co-occurrence KG (baseline).} Add an edge if a predator and prey are mentioned within the same sentence (optionally paragraph). Expected: high coverage, more noise.
  \item \textbf{Method B: Dependency-based relation KG (syntax).} Use dependency parsing and relation patterns (e.g., \textit{``hunts''}, \textit{``preys on''}, \textit{``feeds on''}, including passive forms) to extract explicit predator$\rightarrow$prey edges. Expected: cleaner edges, but may miss relations.
  \item \textbf{Method C: Wikilinks KG (Wikipedia link structure).} Create edges based on Wikipedia internal links between pages (predator page links to prey page). Expected: captures association; not always literal predation.
\end{itemize}

\section*{6. Gold Standard}
To compute Recall and F1, we create a small gold standard of predator$\rightarrow$prey facts, manually extracted from explicit Wikipedia sentences and stored with citations (page URL + sentence snippet). This keeps evaluation reproducible while remaining lightweight.

\section*{7. Evaluation Plan \& Metrics}
For each method, we extract a set of candidate edges and compare them to the gold standard using metrics like Precision, Recall and F1.

To ensure fair comparison, we evaluate only edges where the source is in the predator list and the target is in the prey list.

We also report supporting statistics:
\begin{itemize}[leftmargin=*,nosep]
  \item Runtime per method (seconds) on the same dataset
  \item Number of candidate edges produced per method
  \item Basic descriptive graph properties (e.g., density, average degree)
\end{itemize}

\section*{8. Deliverable}
A single notebook (\texttt{.ipynb}) containing:
\begin{itemize}[leftmargin=*,nosep]
  \item Wikipedia data collection for the selected pages
  \item Implementation of Methods A/B/C
  \item Gold-standard dataset
  \item Precision/Recall/F1 evaluation and summary tables/plots
\end{itemize}

\end{document}
